\section{Conclusion}
\label{sec:conclusion}

QDCVR turns ``the retriever found it'' into ``the system verified it'':
a platform where organization, retrieval, and agency all terminate at an
explicit content-verification gate, where three interchangeable lanes trade
latency for recall under enforced budgets, and where the system's refusal
to speculate is a first-class outcome. Our four research questions are
answered by measurement rather than assertion: verification can be
\emph{enforced} as a platform primitive (RQ1); at the ranking layer the gate
costs well under a second and buys per-citation provenance, but as an ordering
it inherits a same-domain blind spot that we quantify and partially repair
(RQ2); that failure is structural --- it tracks confuser density, not the
checkpoint, domain, language, or scale (RQ3); and honest abstention lives in
the reading--answer layer, never in the threshold, while end-to-end latency
belongs to LLM turns (RQ4).

The immediate roadmap is, in priority order: the missing agent-layer control
--- a grounded lane with the gate disabled (Section~\ref{sec:threats}) --- to
separate retrieval from verification end to end; verifier calibration against
per-library null distributions; trained
evidence-verification models and learned ensembles (the generic-NLI
replay and its span-level, ensemble, and symbolic probes
(Section~\ref{sec:robustness}) already answer the naive upgrades
negatively); combining two-stage
candidates with the verdict-as-filter role that the tool layer suggests;
larger-corpus evaluation with the shipped benchmark suite and an
author-independent question set; consumer-side evidence that
not-found reports serve their readers; and lane auto-selection from
query-shape features --- steps toward retrieval systems whose trustworthiness
is an inspectable property of the pipeline rather than a hope about the
model.

